Aby badać zależność będącą tematem niniejszej pracy, należy w pierwszej kolejności opisać, czym sprzężenie w
ogóle jest, jaka jest jego definicja, a także porównać istniejące definicje w literaturze, ostatecznie formułując
ścisłe miary potrzebne do porównywania alternatywnych wersji programów.
Tym zagadnieniom został poświęcony poniższy rozdział.
Przede wszystkim, stanowi on wprowadzenie do pojęcia sprzężenia pomiędzy klasami.
Przegląda relacje występujące w literaturze i bada ich użyteczność, po czym definiuje intuicję sprzężenia pomiędzy
klasami, wskazując formalne jej definicje.
Stawia tezę, że sprzężenie między klasami może mieć negatywny wpływ na jakość kodu, popierając istniejącymi
badaniami w kontekście obciążenia poznawczego.
W ostatniej części przedstawia dostępne metryki sprzężenia, wskazując ich użyteczne aspekty, a także
przykłady, które uwypuklają różnice między nimi.
Porównano powyższe miary pod kątem ich zastosowania do analizy systemów klas.